\chapter{Συμπεράσματα και Προτάσεις Βελτίωσης}
\label{ch:conclusions}


\section{Προσωρινά συμπεράσματα}
Η ηλεκτρική απόκριση του εδάφους μεταβάλλεται σημαντικά τόσο με τη συχνότητα όσο και με την υγρασία. Στις πολύ χαμηλές συχνότητες κυριαρχεί η πόλωση των ηλεκτροδίων και η ιοντική αγωγιμότητα, ενώ σε υψηλότερες συχνότητες η επίδραση αυτή μειώνεται και αναδεικνύεται περισσότερο η διηλεκτρική συμπεριφορά του νερού.

Τα διαθέσιμα διαγράμματα δείχνουν ότι η αύξηση της υγρασίας μειώνει τις πραγματικές και φανταστικές συνιστώσες της εμπέδησης. Η προσθήκη νερού διευκολύνει την κίνηση διαλυμένων ιόντων και αυξάνει την αγωγιμότητα, με αποτέλεσμα η απόκριση να γίνεται λιγότερο έντονα εξαρτώμενη από τη συχνότητα.

Η αρχική ηλεκτρική μοντελοποίηση υποδεικνύει ότι κανένα απλό ισοδύναμο κύκλωμα δεν περιγράφει κατ' ανάγκη ολόκληρο το εξεταζόμενο φάσμα. Η επιλογή μοντέλου πρέπει να συνδέεται με συγκεκριμένη ζώνη συχνοτήτων και να βασίζεται σε σαφώς ορισμένο κριτήριο προσαρμογής.

\section{Προτάσεις για την ολοκλήρωση της εργασίας}
Για την τελική έκδοση προτείνονται:
\begin{itemize}
  \item πλήρης τεκμηρίωση της προετοιμασίας των δειγμάτων και της γεωμετρίας μέτρησης,
  \item επαναλαμβανόμενες μετρήσεις και παρουσίαση αβεβαιοτήτων,
  \item διαχωρισμός των συχνοτικών περιοχών όπου κυριαρχούν η πόλωση ηλεκτροδίων, η αγωγιμότητα και η διηλεκτρική πόλωση,
  \item προσαρμογή των ισοδύναμων μοντέλων στα πρωτογενή δεδομένα και στατιστική σύγκριση,
  \item βαθμονόμηση της εκτιμώμενης διαπερατότητας ή χωρητικότητας ως προς βαρυμετρική περιεκτικότητα σε νερό και ογκομετρική περιεκτικότητα σε νερό,
  \item επαλήθευση με ανεξάρτητο σύνολο δειγμάτων και, εφόσον είναι δυνατό, με περισσότερους τύπους εδάφους.
\end{itemize}

\section{Επίλογος}
Η φασματοσκοπία εμπέδησης αποτελεί υποσχόμενη προσέγγιση για χαμηλού κόστους εκτίμηση της υγρασίας, αλλά η αξιοπιστία της εξαρτάται από τη βαθμονόμηση, τη συχνότητα λειτουργίας, τη γεωμετρία των ηλεκτροδίων και τις φυσικοχημικές ιδιότητες του εδάφους.
